\documentclass[11pt,a4paper]{article}

\usepackage[T1]{fontenc}
\usepackage[utf8]{inputenc}
\usepackage{lmodern}
\usepackage{microtype}
\usepackage[a4paper,margin=24mm,headheight=15pt]{geometry}
\usepackage{xcolor}
\usepackage{amsmath,amssymb,mathtools}
\usepackage{booktabs,longtable,tabularx,array,multirow}
\usepackage{enumitem}
\usepackage{fancyhdr}
\usepackage{hyperref}
\usepackage{url}
\usepackage{listings}
\usepackage[most]{tcolorbox}
\usepackage{titlesec}
\usepackage{lastpage}
\usepackage{pifont}
\usepackage{seqsplit}

\definecolor{AionNavy}{HTML}{13233A}
\definecolor{AionBlue}{HTML}{245A8D}
\definecolor{AionTeal}{HTML}{168C8C}
\definecolor{AionGreen}{HTML}{237A57}
\definecolor{AionAmber}{HTML}{B97618}
\definecolor{AionRed}{HTML}{A63D40}
\definecolor{AionInk}{HTML}{1F2933}
\definecolor{AionGray}{HTML}{667788}
\definecolor{AionPale}{HTML}{F3F7FA}
\definecolor{AionRule}{HTML}{D8E2EA}

\hypersetup{
  colorlinks=true,
  linkcolor=AionBlue,
  urlcolor=AionTeal,
  citecolor=AionBlue,
  pdftitle={AION Riemann Hypothesis Programme: Research Record and Roadmap},
  pdfauthor={AION / COMDEX},
  pdfsubject={Formal mathematics research programme status, evidence, and roadmap},
  pdfkeywords={AION, Riemann Hypothesis, Lean 4, Mathlib, formal verification, AI mathematics}
}

\pagestyle{fancy}
\fancyhf{}
\lhead{\textcolor{AionGray}{AION Riemann Hypothesis Programme}}
\rhead{\textcolor{AionGray}{Research Record - 15 August 2026}}
\cfoot{\textcolor{AionGray}{\thepage\ / \pageref{LastPage}}}
\renewcommand{\headrulewidth}{0.4pt}
\renewcommand{\headrule}{\hbox to\headwidth{\color{AionRule}\leaders\hrule height \headrulewidth\hfill}}

\titleformat{\section}{\Large\bfseries\color{AionNavy}}{\thesection}{0.7em}{}
\titleformat{\subsection}{\large\bfseries\color{AionBlue}}{\thesubsection}{0.7em}{}
\titleformat{\subsubsection}{\normalsize\bfseries\color{AionInk}}{\thesubsubsection}{0.7em}{}
\titlespacing*{\section}{0pt}{2.3ex plus .4ex}{1.0ex}
\titlespacing*{\subsection}{0pt}{1.8ex plus .3ex}{0.7ex}

\setlist[itemize]{leftmargin=1.5em,itemsep=0.25em,topsep=0.4em}
\setlist[enumerate]{leftmargin=1.7em,itemsep=0.35em,topsep=0.4em}
\renewcommand{\arraystretch}{1.18}
\setlength{\tabcolsep}{2pt}
\setlength{\parindent}{0pt}
\setlength{\parskip}{0.55em}

\newcommand{\pass}{\textcolor{AionGreen}{\ding{51}}}
\newcommand{\fail}{\textcolor{AionRed}{\ding{55}}}
\newcommand{\pending}{\textcolor{AionAmber}{\raisebox{0.15ex}{\rule{0.8ex}{0.8ex}}}}
\newcommand{\code}[1]{\texttt{#1}}
\newcommand{\hash}[1]{\texttt{\footnotesize\seqsplit{#1}}}
\newcommand{\RH}{\mathrm{RH}}
\newcommand{\RePart}{\operatorname{Re}}
\newcommand{\status}[2]{\textbf{#1:} #2}

\newtcolorbox{truthbox}{
  colback=AionRed!4,colframe=AionRed!70!black,
  title=Truth boundary,fonttitle=\bfseries,
  boxrule=0.7pt,arc=2pt,left=6pt,right=6pt,top=5pt,bottom=5pt
}
\newtcolorbox{resultbox}{
  colback=AionGreen!5,colframe=AionGreen!75!black,
  title=Verified result,fonttitle=\bfseries,
  boxrule=0.7pt,arc=2pt,left=6pt,right=6pt,top=5pt,bottom=5pt
}
\newtcolorbox{decisionbox}{
  colback=AionBlue!4,colframe=AionBlue!75!black,
  title=Programme decision,fonttitle=\bfseries,
  boxrule=0.7pt,arc=2pt,left=6pt,right=6pt,top=5pt,bottom=5pt
}
\newtcolorbox{warningbox}{
  colback=AionAmber!6,colframe=AionAmber!80!black,
  title=Open gate,fonttitle=\bfseries,
  boxrule=0.7pt,arc=2pt,left=6pt,right=6pt,top=5pt,bottom=5pt
}

\lstdefinelanguage{Lean}{
  morekeywords={theorem,lemma,def,by,exact,import,namespace,end,have,simpa,using,fun,if,then,else},
  sensitive=true,
  morecomment=[l]{--},
  morestring=[b]"
}
\lstset{
  language=Lean,
  basicstyle=\ttfamily\small,
  keywordstyle=\color{AionBlue}\bfseries,
  commentstyle=\color{AionGray}\itshape,
  stringstyle=\color{AionTeal},
  backgroundcolor=\color{AionPale},
  frame=single,
  rulecolor=\color{AionRule},
  breaklines=true,
  columns=fullflexible,
  keepspaces=true,
  showstringspaces=false,
  xleftmargin=0.5em,xrightmargin=0.5em,
  aboveskip=0.8em,belowskip=0.8em
}

\begin{document}

\begin{titlepage}
  \pagecolor{AionNavy}
  \color{white}
  \vspace*{18mm}
  {\fontsize{23}{28}\selectfont\bfseries AION Riemann Hypothesis Programme\par}
  \vspace{4mm}
  {\fontsize{15.5}{19}\selectfont Formal Mathematics Research Record, Evidence Review, and Roadmap\par}
  \vspace{15mm}
  {\color{AionTeal!45!white}\rule{\textwidth}{1.4pt}}
  \vspace{10mm}

  \begin{tabularx}{\textwidth}{>{\bfseries}p{38mm}X}
    Document date & 15 August 2026 \\
    Programme & AION formal Riemann Hypothesis research \\
    Repository & \path{/Users/kevinrobinson/dev/COMDEX} \\
    Formal authority & Lean 4.24.0 kernel and Mathlib v4.32.1 \\
    Current stage & Compact zeta divisor complete; explicit-formula bridge open \\
    Research status & No new RH mathematics; RH remains unsolved \\
  \end{tabularx}

  \vfill
  \begin{tcolorbox}[colback=AionBlue!55!AionNavy,colframe=AionTeal!70!white,title=Purpose of this record,
    fonttitle=\bfseries,coltitle=white,coltext=white,boxrule=0.6pt]
  This document records what AION has actually built, measured, learned, and not
  established. It is designed to survive handovers, prevent claim inflation,
  support reproducibility, and define the shortest credible route from formal
  proof competence to useful research on a genuine RH prerequisite.
  \end{tcolorbox}
  \vspace{5mm}
  {\small Internal research record - evidence-first, kernel-governed}
\end{titlepage}
\nopagecolor
\color{AionInk}

\pagenumbering{roman}
\tableofcontents
\clearpage

\section*{Document control}
\addcontentsline{toc}{section}{Document control}

\begin{tabularx}{\textwidth}{>{\bfseries}p{37mm}X}
Version & 1.3 \\
Status & Full Chebyshev Perron formula kernel-checked; contour orientation and the residues at nontrivial zeros, $s=1$, and $s=0$ checked \\
Scope & Formal RH programme only; excludes unrelated prime-search, factoring, visual probes, and Photon Algebra work \\
Primary evidence & Lean sources, pinned environment, immutable JSON results, SHA-256 hashes, focused tests \\
Interpretive authority & The Lean kernel decides proof acceptance; this report interprets experiments but cannot create mathematical truth \\
Change rule & Future revisions must preserve prior immutable evidence and clearly separate new measurements from retrospective interpretation \\
\end{tabularx}

\begin{truthbox}
AION has not solved the Riemann Hypothesis. No new theorem that advances a
resolution of RH has been established. Reconstructing Mathlib theorems is
evidence about proof-system competence and experimental plumbing, not original
mathematics. Numerical evidence, model confidence, persuasive prose, and
self-awarded scores are not proofs.
\end{truthbox}

\clearpage
\pagenumbering{arabic}

\section{Executive summary}

The AION RH programme now has a credible formal laboratory, a pinned Lean and
Mathlib authority path, strict rejection of incomplete proofs, immutable
matched experiments, governed retrieval from AION's persistent knowledge
store, and an extended kernel-checked mathematics chain reaching from zeta-zero
infrastructure to finite Chebyshev--Perron inversion.

The strongest verified engineering result is not a theorem about RH. It is a
research-control result: the system can now distinguish a proposed Lean proof
from an accepted proof, retain only compiler-verified method metadata, expose
exact provenance identifiers to experimental arms, and fail closed before a
paid run if task or environment hashes do not match.

The one paid experiment completed on 14 August 2026 did not establish an AION
advantage. The proposer-only arm solved 0 of 3 tasks. Each AION-labelled arm
solved 1 of 3. Those AION labels were not yet causally clean: memory was a
generic prompt hint, repair was ordinary compiler feedback, and one task failed
because the generated source omitted the Mathlib module containing the needed
theorem. The result therefore validates the live provider-to-Lean path but does
not measure the intended architecture.

The continuation work corrected that defect and completed both the v3 matched
comparison and its delayed source-closed retention test. Neither passed its
precommitted superiority gate. These negative results are preserved, and the
programme has moved to the mathematical dependency chain without converting an
engineering signal into a capability claim.

The strongest new formal result is now the complete Chebyshev Perron formula.
Lean verifies the squared-Cauchy Fourier value, the sharp away-from-jump scalar
limit, finite Dirichlet-polynomial inversion, fixed-height interchange of the
full von Mangoldt series and vertical integral, a height-independent summable
majorant, and the limit to $\psi(x)$ for nonintegral $x>1$. Initial contour
infrastructure is also checked: exact rectangle orientation, identification of
the right edge with the proved Perron integral, multiplicity-aware zero
residues, the main-term residue $x$ at $s=1$, and the $-\log(2\pi)$ residue at
$s=0$.

\begin{resultbox}
Final formal verification status: \textbf{PASS}. The complete RH library builds
in the pinned environment in 3,590 jobs. The Fourier, full Perron, compact
divisor, and initial contour-residue results contain no \code{sorry}, \code{admit},
project axiom, unsafe declaration, paid call, or hidden oracle step.
\end{resultbox}

\subsection{Status at a glance}

\begin{tabularx}{\textwidth}{X c p{55mm}}
\toprule
Capability or claim & Status & Interpretation \\
\midrule
Pinned formal laboratory & \pass & Lean and Mathlib are operational authorities \\
Strict no-placeholder policy & \pass & \code{sorry}, \code{admit}, and local axioms are prohibited \\
Known zeta baseline compiles & \pass & Established mathematics is reconstructed \\
Live provider-to-Lean experiment & \pass & One paid run completed and evidence is immutable \\
Genuine verified-memory retrieval & \pass & AION persistent knowledge retrieval is now wired \\
Bounded diagnostic repair records & \pass & Failure classes produce auditable routes \\
Six-family v2 packet locally verified & \pass & All hidden oracle proofs compile locally \\
V2/v3 governed execution & \pass & Sealed comparison completed with immutable evidence \\
Delayed source-closed retention & \pass & Completed once; memory arm scored 2/6 versus cold 3/6 \\
Compact multiplicity-aware zeta divisor & \pass & Finite compact zero sums are kernel-checked \\
Scalar Perron theorem & \pass & Sharp limit proved away from $y=1$ \\
Finite Chebyshev--Perron inversion & \pass & Finite von Mangoldt specialization converges to $\psi(x)$ \\
Full Chebyshev Perron formula & \pass & Growing-height interchange and uniform domination proved \\
Multiplicity-aware local zero residues & \pass & Residue is $-m(\rho)x^\rho/\rho$ \\
Pole residues at $s=1$ and $s=0$ & \pass & Contributions $x$ and $-\log(2\pi)$ checked \\
Global contour shift and explicit formula & \pending & Pole excision, trivial-zero sum, and edge limits remain \\
AION superiority established & \fail & Neither precommitted experiment met the gate \\
New original RH mathematics & \fail & Classical dependencies formalized; no novel theorem \\
RH solved & \fail & No \\
\bottomrule
\end{tabularx}

\section{Mathematical target and scope}

\subsection{The Riemann zeta function}

For complex $s$ with $\RePart(s)>1$, the Riemann zeta function is represented by
the absolutely convergent Dirichlet series
\[
  \zeta(s)=\sum_{n=1}^{\infty}\frac{1}{n^s}.
\]
It admits a meromorphic continuation to the complex plane, with a simple pole
at $s=1$. Its nontrivial zeros lie in the critical strip
\[
  0 < \RePart(s) < 1.
\]
The Riemann Hypothesis asserts
\[
  \RH:\qquad \zeta(s)=0,\ 0<\RePart(s)<1
  \quad\Longrightarrow\quad
  \RePart(s)=\frac12.
\]

This project does not treat the displayed implication as a prompt to be solved
directly. The target is too deep, existing formal-library coverage is partial,
and an unconstrained language model can easily produce plausible but invalid
mathematics. The programme therefore decomposes the challenge into governed,
machine-checkable competence stages.

\subsection{What counts as progress}

Four evidence classes must not be conflated:

\begin{enumerate}
  \item \textbf{Infrastructure progress}: the toolchain, verifier, hashing,
  memory, repair, and experiment controls work correctly.
  \item \textbf{Competence progress}: AION reconstructs or transfers established
  proofs more reliably, cheaply, or durably than matched controls.
  \item \textbf{Formalization progress}: AION formalizes known mathematics that
  was missing from the local dependency chain. This can be useful and novel as
  engineering without being new mathematics.
  \item \textbf{Mathematical progress}: a genuinely new lemma, method, bound, or
  theorem relevant to RH survives Lean verification and independent expert
  mathematical review.
\end{enumerate}

Current achievement is primarily in classes 1 and 2's experimental
prerequisites. There is not yet evidence in class 4.

\section{Programme thesis}

The programme's central hypothesis is architectural:

\begin{quote}
Can AION turn formally verified mathematical experience into retained,
repairable competence under source closure, beyond what the same underlying
proposer achieves alone?
\end{quote}

The governed loop is
\[
\small
\begin{aligned}
\boxed{\text{task}}&\rightarrow\boxed{\text{proposal}}
\rightarrow\boxed{\text{Lean gate}}
\rightarrow\boxed{\text{diagnosis}}
\rightarrow\boxed{\text{bounded repair}}
\rightarrow\boxed{\text{verified memory}}\\[-1mm]
&\hspace{55mm}\Downarrow\\[-1mm]
&\hspace{42mm}\boxed{\text{source-closed transfer}}
\end{aligned}
\]

The same proposer, model version, temperature, task, tool access, token budget,
attempt budget, and time budget must be held fixed. Only AION memory and repair
mechanisms may vary. This is what makes an observed difference attributable to
the architecture rather than to a stronger model or an easier prompt.

\subsection{Progression of research difficulty}

\begin{enumerate}
  \item Known zeta and complex-analysis micro-lemmas.
  \item Held-out reformulations and theorem specialization.
  \item Delayed reconstruction with training sources closed.
  \item Formal RH-equivalent statements and dependency mapping.
  \item Complete reconstruction of a selected prerequisite chain.
  \item Identification of one precise missing formalization or mathematical
  bottleneck.
  \item Formalization and independent review of useful supporting lemmas.
  \item Only then, carefully governed work on a genuinely open bottleneck.
\end{enumerate}

\section{Trusted formal laboratory}

\subsection{Authority model}

The declared trusted base is the pinned Lean kernel and pinned Mathlib
dependency. A candidate counts only when the Lean process returns success for
the complete generated source. The harness also retains a SHA-256 source hash
and compiler diagnostic.

Project-authored proof sources are scanned for prohibited incompleteness:
\begin{itemize}
  \item \code{sorry};
  \item \code{admit};
  \item project-local \code{axiom} declarations;
  \item hidden placeholders or self-awarded acceptance.
\end{itemize}

This policy intentionally does not scan upstream Mathlib as though its trusted
implementation were AION-authored source. The kernel return code remains the
final acceptance signal.

\subsection{Established baseline}

The baseline reconstructs three established facts already present in Mathlib:

\begin{align*}
  \zeta(0) &= -\frac12,\\
  \zeta\!\left(-2(n+1)\right) &= 0 \qquad (n\in\mathbb{N}),\\
  \RePart(s)>1 &\Longrightarrow \zeta(s)\neq 0.
\end{align*}

Representative Lean source:

\begin{lstlisting}
theorem reproduce_zeta_at_zero : riemannZeta 0 = -(1 : C) / 2 := by
  exact riemannZeta_zero

theorem reproduce_trivial_zero (n : N) :
    riemannZeta (-2 * ((n : C) + 1)) = 0 := by
  exact riemannZeta_neg_two_mul_nat_add_one n

theorem reproduce_zero_free_right_half_plane {s : C} (hs : 1 < s.re) :
    riemannZeta s != 0 := by
  exact riemannZeta_ne_zero_of_one_lt_re hs
\end{lstlisting}

In the actual Lean files, \code{C}, \code{N}, and \code{!=} above are the
corresponding Unicode complex-number, natural-number, and not-equal symbols.
They are rendered here in ASCII for portable documentation.

\subsection{Fixture smoke harness}

The fixture-only harness validates Lean compilation, proof-policy enforcement,
evidence hashing, source-closed fixture reconstruction, repair recording, and
claim boundaries. Its status is \code{PASS\_TRUSTED\_LAB\_SMOKE\_TEST}.

This is deliberately a plumbing test. Fixed candidates do not measure model
quality, AION intelligence, memory advantage, repair advantage, transfer, or
retention.

\section{First live matched experiment}

\subsection{Design}

One live experiment used GPT-4.1 through the OpenAI Responses API. Four arms
received the same three established theorem tasks, with at most two attempts per
task:

\begin{enumerate}
  \item proposer alone;
  \item AION-labelled memory plus repair;
  \item AION-labelled repair without memory;
  \item AION-labelled memory without repair.
\end{enumerate}

The proposer returned proof bodies. The harness generated complete Lean files,
rejected prohibited tokens, invoked the pinned compiler, and retained prompts,
proofs, diagnostics, hashes, provider identifiers, token usage, and elapsed
time.

\subsection{Measured result}

\begin{table}[ht]
\centering
\caption{Paid experiment result, 14 August 2026}
\begin{tabular}{lrrrr}
\toprule
Arm & Accepted & Tasks & Model calls & Success rate \\
\midrule
Proposer alone & 0 & 3 & 6 & 0\% \\
AION memory + repair label & 1 & 3 & 5 & 33.3\% \\
AION without memory label & 1 & 3 & 5 & 33.3\% \\
AION without repair label & 1 & 3 & 5 & 33.3\% \\
\bottomrule
\end{tabular}
\end{table}

Additional measurements:
\begin{itemize}
  \item resolved provider model: \code{gpt-4.1-2025-04-14};
  \item total provider tokens: 5,449;
  \item recorded wall-clock duration: 134.46 seconds;
  \item accepted family: only the elementary value at zero;
  \item immutable result hash:
  \hash{5795e3faa55cf2a7466a3537ac41005056d44711eaa826c84935af318437b2e6}.
\end{itemize}

\begin{decisionbox}
The experiment does not support an AION superiority claim. The three
AION-labelled arms tied, and the labels did not yet correspond to the genuine
durable-memory and governed-repair mechanisms. The correct conclusion is a
null architectural result with useful infrastructure evidence.
\end{decisionbox}

\subsection{Failure taxonomy}

Rejected candidates exhibited several recurring failure modes:

\begin{itemize}
  \item invented or misremembered theorem identifiers;
  \item Lean tactic syntax errors and invalid indentation;
  \item incorrect binders or explicit arguments;
  \item theorem conclusions that did not match the generated goal;
  \item unsolved goals after partial rewriting;
  \item repeated use of a failed branch instead of theorem discovery.
\end{itemize}

These failures show why a generic ``try again'' loop is inadequate. The system
needs catalogue discovery, exact signature inspection, diagnostic
classification, branch abandonment, and verified memory.

\section{Root-cause correction and continuation achievements}

\subsection{Concrete environment defect}

The local baseline compiled, but the live generator imported only the
\code{RiemannZeta} module. The theorem
\code{riemannZeta\_ne\_zero\_of\_one\_lt\_re} is made available through the
Mathlib \code{Dirichlet} module. The exact live right-half-plane template
therefore failed even with a known-correct proof term.

The generator now imports both required modules, and regression tests compile
known-correct proofs through the exact live templates. This finding changes the
interpretation of the earlier failure: at least one failure attributed to the
proposer was actually an experiment-environment error.

\subsection{Truthful experiment state}

The original packet incorrectly remained marked \code{not\_run} after the paid
result existed. It now records the completed run, its timestamped evidence
path, original sealed-packet hash, result hash, model, timestamps, and a
non-superiority interpretation. The immutable result itself was not rewritten.

\subsection{Genuine governed memory retrieval}

A scoped adapter now uses AION's
\code{HexCorePersistentLearningRuntime.knowledge.retrieve} mechanism. Admission
is fail closed:

\begin{enumerate}
  \item the exact theorem template is compiled by Lean;
  \item the return code must be zero;
  \item a 64-character source SHA-256 hash must be present;
  \item only then may a verified evidence capsule and claim be ingested;
  \item proof bodies are not retained in the memory object;
  \item retrieval returns claim IDs, capsule IDs, evidence checksums, method
  metadata, and the verified source hash.
\end{enumerate}

The clean scoped store currently contains:

\begin{center}
\begin{tabular}{lr}
\toprule
Metric & Value \\
\midrule
Active verified claims & 6 \\
Verified evidence capsules & 6 \\
Contradictions & 0 \\
Store revision after clean rebuild & 6 \\
Retained proof bodies & 0 \\
\bottomrule
\end{tabular}
\end{center}

The six acquisition families cover the zeta value at zero, trivial zeros,
right-half-plane non-vanishing, differentiability away from one,
right-half-plane positivity, and finiteness of zeta zeros in compact sets.

\subsection{Governed Lean repair routing}

The repair controller classifies diagnostics into explicit routes:

\begin{longtable}{p{35mm}p{50mm}p{66mm}}
\toprule
Category & Detection signal & Required response \\
\midrule
\endhead
Unknown theorem & Unknown identifier diagnostic & Search installed Mathlib, inspect exact signature, replace invented names \\
Syntax failure & Unexpected token, invalid tactic, malformed command & Correct Lean syntax and return one well-indented tactic body \\
Type mismatch & Application mismatch, failed synthesis, coercion problem & Inspect goal and theorem types; repair arguments and coercions \\
Unsolved goal & Remaining goals after a partial proof & Split the goal or apply a verified theorem with a matching conclusion \\
Unclassified or repeated failure & No trusted route or recurring branch & Abandon the branch and construct a materially different minimal proof \\
\bottomrule
\end{longtable}

Each repair record contains a stable record ID, task ID, attempt number,
category, diagnostic hash, instruction, and a one-attempt retry bound. A
repair-enabled arm receives the record and compiler diagnostic; a no-repair arm
does not.

\subsection{Fail-closed paid execution}

The live harness now refuses to load an API key unless all of the following are
true:

\begin{itemize}
  \item an explicitly selected packet is marked \code{ready};
  \item every task is source closed;
  \item every statement hash matches;
  \item the pinned Lean environment hash matches;
  \item the verified AION memory store is nonempty;
  \item the packet contains normalized tasks and matched arms.
\end{itemize}

This turns ``do not run yet'' from a procedural note into an executable safety
property.

\section{V2 held-out experiment}

\subsection{Purpose}

The v2 packet replaces the three direct reconstruction tasks with six
reformulations across six families. Each statement has been independently
compiled with a known-correct oracle proof, but the proposer does not receive
the proof body. Statements and environment are hash frozen.

\begin{longtable}{@{}p{10mm}p{48mm}p{91mm}@{}}
\toprule
Ref. & Family & Held-out statement \\
\midrule
\endhead
T1 & Special-value transfer & $2\zeta(0)=-1$ \\
T2 & Trivial-zero specialization & $\zeta(-6)=0$ \\
T3 & Non-vanishing contrapositive & $\zeta(s)=0\Rightarrow\RePart(s)\leq1$ \\
T4 & Analytic prerequisite & $s\neq1\Rightarrow\zeta$ differentiable at $s$ \\
T5 & Right-half-plane positivity & $x>1\Rightarrow\RePart(\zeta(x))>0$ \\
T6 & Zero-set topology & Compact sets meet the zeta zero set in finitely many points \\
\bottomrule
\end{longtable}

\textbf{Packet IDs:}
\begin{itemize}[leftmargin=11mm,itemsep=0pt,topsep=2pt]
  \item[T1] \code{scaled\_zeta\_at\_zero}
  \item[T2] \code{trivial\_zero\_at\_neg\_six}
  \item[T3] \code{zero\_implies\_re\_le\_one}
  \item[T4] \code{zeta\_differentiable\_away\_from\_one}
  \item[T5] \code{zeta\_real\_part\_positive}
  \item[T6] \code{compact\_has\_finitely\_many\_zeta\_zeros}
\end{itemize}

\subsection{Matched arms}

The precommitted packet defines:

\begin{enumerate}
  \item proposer alone;
  \item full AION memory plus repair;
  \item AION repair without memory;
  \item AION memory without repair;
  \item a cold, isolated-state control.
\end{enumerate}

Each arm receives at most two attempts per task with equal token, time, and tool
budgets. Exam-time memory writes are disabled, preventing early arms or tasks
from leaking answers into later conditions.

\subsection{Precommitted scoring}

The primary score is the number of Lean-accepted tasks per arm. Secondary
measures are attempts, provider tokens, and elapsed time. A superiority claim
requires the full AION arm to exceed \emph{every} matched ablation on the
primary score. A tie is not evidence of superiority.

The packet remains \code{local\_verification\_only}. It must not be changed to
\code{ready} until an independent sealing review confirms:

\begin{itemize}
  \item no oracle proof or answer text reaches the proposer;
  \item task and environment hashes match the reviewed packet;
  \item model and sampling parameters are frozen;
  \item arm order and potential order effects are addressed;
  \item the delayed-retention schedule is precommitted;
  \item output locations and immutable evidence rules are confirmed.
\end{itemize}

\section{Evaluation methodology}

\subsection{Causal questions}

The next experiment should answer four distinct questions:

\begin{description}[leftmargin=37mm,style=nextline]
  \item[Reconstruction] Can the proposer produce kernel-accepted established proofs?
  \item[Memory effect] Does verified retrieval improve success or efficiency over an identical no-memory arm?
  \item[Repair effect] Does diagnostic routing improve rejected candidates over an identical no-repair arm?
  \item[Retention effect] Does the method survive source closure and time delay on equivalent reformulations?
\end{description}

Immediate success without delayed retention is not sufficient evidence for the
programme's main thesis. The decisive property is retained, source-closed
competence.

\subsection{Required measurements}

For every arm, task, and attempt, retain:

\begin{itemize}
  \item task, environment, prompt, source, and output hashes;
  \item model alias and resolved provider model version;
  \item temperature, seed if supported, token cap, and tool policy;
  \item exact memory claim IDs, capsule IDs, and evidence checksums supplied;
  \item exact repair record ID, category, diagnostic hash, and route supplied;
  \item candidate proof body and complete bounded compiler diagnostic;
  \item Lean return code and generated source hash;
  \item provider tokens, model calls, elapsed time, and any retries;
  \item failure classification and branch-abandonment decision;
  \item immutable result path, final result hash, and prior-result linkage.
\end{itemize}

\subsection{Statistical limitations}

Six tasks are still a small sample. Even a perfect 6--0 result would not prove
broad mathematical intelligence. The first v2 run should be treated as a
bounded mechanistic study. If promising, it should be replicated on a larger,
pre-registered curriculum with multiple theorem families, multiple random task
draws, and delayed tests. Confidence intervals and paired task-level analyses
should replace informal percentage comparisons as the sample grows.

\section{Roadmap toward meaningful RH work}

\subsection{Stage gates}

\begin{longtable}{p{12mm}p{42mm}p{72mm}p{28mm}}
\toprule
Stage & Objective & Exit criterion & Current \\
\midrule
\endhead
0 & Trusted formal laboratory & Pinned compiler, strict policy, immutable evidence & \pass \\
1 & Exact-template reconstruction & Known-correct terms compile through live generator & \pass \\
2 & Real memory and repair & Provenance IDs and repair records differ causally by arm & \pass \\
3 & Sealed immediate transfer & Run and score the precommitted held-out comparison & \pass \\
4 & Delayed retention & Fresh Lean acceptance after source closure and delay & \pass \\
5 & Critical-strip dependency reconstruction & Formal map through compact divisors and Perron interfaces & \pass \\
6 & Missing formalization & Scalar Perron and finite Chebyshev inversion kernel-checked & \pass \\
7 & Genuine bottleneck lemma & New candidate survives Lean and expert review & \pending \\
8 & RH-level contribution & Independent community validation of material progress & \pending \\
\bottomrule
\end{longtable}

\subsection{Immediate next steps}

\textbf{Phase A - Complete the Chebyshev Perron bridge (completed)}
\begin{enumerate}
  \item Prove fixed-height interchange of the absolutely convergent von
  Mangoldt coefficient series and the vertical interval integral.
  \item Establish an explicit coefficient-tail bound uniform in the truncation
  height, or an equivalent two-parameter convergence theorem.
  \item Combine that bound with finite Chebyshev inversion to prove
  \code{ChebyshevPerronInversionTarget} for nonintegral $x>1$ and $c>1$.
\end{enumerate}

\textbf{Phase B - Shift the contour and recover residues (active)}
\begin{enumerate}
  \item Construct admissible rectangles avoiding poles and nontrivial zeros.
  \item Prove the horizontal and displaced vertical boundary estimates.
  \item Use the checked residues at $s=1$, at $s=0$, and at every nontrivial
  zero; prove simplicity and sum the negative-even trivial-zero residues.
  \item Pass from finite compact zero sums to the selected explicit-formula
  convergence convention.
\end{enumerate}

\textbf{Phase C - Reach the selected RH-equivalent criterion}
\begin{enumerate}
  \item Derive the Chebyshev $\psi$ remainder from the explicit formula.
  \item Formalize the RH-to-square-root-remainder direction.
  \item Formalize the converse implication from the remainder bound to
  localization of every nontrivial zero on $\RePart(s)=1/2$.
  \item Submit every new mathematical lemma to independent expert review before
  raising its claim level.
\end{enumerate}

\subsection{Longer-term research directions}

If the competence gates pass, useful research directions include:

\begin{itemize}
  \item formal equivalences of RH and careful comparison of their dependency
  costs;
  \item explicit zero symmetry and critical-strip localization lemmas;
  \item formal bridges between Dirichlet-series, Euler-product, and analytic
  continuation representations;
  \item reusable bounds needed for zero-free regions or contour arguments;
  \item formalization of classical lemmas absent from Mathlib;
  \item counterexample search for proposed intermediate lemmas;
  \item certified numerical bounds using Arb/FLINT as supporting evidence,
  never as a substitute for proof;
  \item automated theorem discovery over installed Mathlib declarations and
  exact signature retrieval;
  \item proof-term minimization and robustness testing under import reduction.
\end{itemize}

\section{Research governance}

\subsection{Claim ladder}

Public and internal statements should use the narrowest justified claim:

\begin{longtable}{p{45mm}p{30mm}p{80mm}}
\toprule
Claim & Allowed now? & Required evidence \\
\midrule
\endhead
The trusted RH lab compiles known results & Yes & Pinned Lean acceptance and source hashes \\
A live provider-to-Lean run completed & Yes & Immutable timestamped result and provider records \\
AION memory and repair are wired & Yes, narrowly & Code, local verification, provenance logs, ablation design \\
AION improves this proposer & No & Full arm exceeds every matched ablation under precommitment \\
AION retains mathematical methods & No & Delayed source-closed reconstruction \\
AION formalized a useful known lemma & No, not yet & New Lean source, dependency impact, independent review \\
AION discovered new mathematics & No & Novelty analysis, Lean proof, expert review, literature search \\
AION advanced RH & No & Material relevance to a recognized bottleneck and external validation \\
AION solved RH & No & Complete proof, formal verification, extensive independent mathematical review \\
\bottomrule
\end{longtable}

\subsection{Risk register}

\begin{longtable}{p{38mm}p{18mm}p{95mm}}
\toprule
Risk & Level & Mitigation \\
\midrule
\endhead
Invented theorem names & High & Installed-library search and exact signature inspection before proposal \\
Import/environment mismatch & High & Compile oracle terms through exact generated templates; freeze environment hash \\
Answer leakage into memory & High & Store verified metadata, not proof bodies; disable exam writes; source-close tasks \\
Arm contamination & High & Isolated state, immutable packets, logged retrieval IDs, fixed order or randomized precommitment \\
Repeated paid reruns & High & Exactly one bounded run; fail-closed status gate; immutable reporting \\
Claim inflation from ties & High & No-tie superiority rule and explicit claim matrix \\
Small-sample overinterpretation & Medium & Treat v2 as mechanistic; expand only after precommitted success \\
Mutable or lost evidence & High & Timestamped immutable results, hashes, version control, append-only lineage \\
Untracked repository work & High & Review and commit RH assets separately from unrelated dirty-tree changes \\
Numerics mislabelled as proof & High & Separate certified-numerics lab and explicit non-proof label \\
Formal proof of irrelevant lemma & Medium & Dependency graph and external relevance review before promotion \\
Self-grading or circular authority & High & Kernel acceptance plus independent mathematical review for research claims \\
\bottomrule
\end{longtable}

\subsection{Operating discipline}

The following rules are part of the scientific method, not merely workflow
preferences:

\begin{itemize}
  \item Never say a test is running unless a real process or session exists.
  \item Never spend API credit before all local gates pass.
  \item Do not repeatedly announce the same next step.
  \item If an edit fails, record the exact failure once, correct it, and verify.
  \item Do not mutate a frozen mission after observing failures.
  \item Do not delete or rewrite immutable results.
  \item Do not turn a tie, failure, or confounded result into a success claim.
  \item Keep consciousness, proof-of-life, and self-awareness claims outside the
  RH evaluation.
  \item Separate known-mathematics reconstruction, new formalization, and new
  mathematics in every report.
\end{itemize}

\section{Repository and evidence map}

\subsection{Core laboratory}

\begin{description}[leftmargin=40mm,style=nextline]
  \item[\code{research/riemann\_lab/README.md}] Scope, proof policy, and claim boundary.
  \item[\code{lean-toolchain}] Pinned Lean toolchain.
  \item[\code{lakefile.toml}] Lean project and Mathlib requirement.
  \item[\code{lake-manifest.json}] Resolved dependency state.
  \item[\code{EstablishedZeta.lean}] Kernel-accepted established baseline.
  \item[\code{CatalogProbe.lean}] Exact theorem-name and signature probes.
  \item[\code{dependency\_graph.json}] RH prerequisite graph and selected critical-strip target.
  \item[\code{live\_task\_packet.json}] Truthful record of the completed v1 run.
  \item[\code{live\_task\_packet\_v3.json}] Completed hash-frozen six-family comparison packet.
\end{description}

\subsection{Harnesses and tests}

\begin{description}[leftmargin=44mm,style=nextline]
  \item[\code{scripts/aion\_riemann\_lab.py}] Fixture-only trusted smoke harness.
  \item[\code{scripts/aion\_riemann\_live\_exam.py}] Live matched experiment and fail-closed preflight.
  \item[\code{scripts/aion\_riemann\_verified\_memory.py}] Governed memory adapter and Lean repair controller.
  \item[\code{scripts/aion\_riemann\_local\_verify.py}] Unpaid authority audit, clean memory build, and report generation.
  \item[\code{backend/tests/test\_aion\_riemann\_lab.py}] Fourteen focused policy and regression tests.
\end{description}

\subsection{Evidence}

\begin{description}[leftmargin=55mm,style=nextline]
  \item[\code{live\_exam\_20260814T213917Z.json}] Immutable paid result.
  \item[\code{live\_exam\_latest.json}] Latest pointer with identical content hash.
  \item[\code{local\_verification\_latest.json}] Final zero-cost local verification evidence.
  \item[\code{aion\_verified\_memory.json}] Scoped six-claim verified memory store.
  \item[\code{AION\_RIEMANN\_LOCAL\_VERIFICATION\_2026-08-15.md}] Human-readable continuation report.
\end{description}

Older visual probes, Photon Algebra drivers, and other historical scripts are
outside the trusted programme and must not be silently merged into this
evidence chain.

\section{Open questions}

The programme should explicitly track the following unresolved questions:

\begin{enumerate}
  \item Does verified memory help when the held-out statement is a nontrivial
  reformulation rather than a direct theorem-name lookup?
  \item Does diagnostic routing add value beyond raw compiler feedback?
  \item How sensitive are results to task order and retrieved-memory ranking?
  \item Should theorem discovery be a deterministic library-search tool rather
  than prompt-visible signature memory, or should both be separate arms?
  \item How should cold AION differ operationally from proposer alone and from
  no-memory AION?
  \item What delay is long enough to make source-closed retention meaningful?
  \item What minimum task count supports a credible architectural claim?
  \item Which critical-strip dependency is both missing from Mathlib and useful
  enough to justify formalization effort?
  \item How will novelty be checked before any claim of new mathematics?
  \item Which external mathematicians or formalization experts will review
  candidate bottleneck lemmas?
\end{enumerate}

\section{Recommended research log schema}

Every future session should append a compact structured record with:

\begin{itemize}
  \item date, operator, repository commit, dirty-tree status, and toolchain hash;
  \item objective and explicit non-objectives;
  \item packet hash, task hashes, arm definitions, model settings, and budget;
  \item process/session identifier for any live execution;
  \item exact commands or harness entry points used;
  \item local test result before paid work;
  \item immutable evidence paths and hashes after work;
  \item measured outcomes, failures, and unresolved anomalies;
  \item claim classification: infrastructure, competence, formalization, or mathematics;
  \item decision: stop, replicate, revise locally, advance, or abandon;
  \item next stage gate and the evidence required to pass it.
\end{itemize}

This log prevents narrative drift and makes negative results scientifically
useful.

\section{Final programme assessment}

The AION RH project has crossed an important but preliminary threshold. It is
no longer an unconstrained request for an AI to speculate about RH. It is a
formal, evidence-bearing programme capable of rejecting invalid proofs,
preserving negative results, recording exact provenance, and testing specific
architectural hypotheses.

The current accomplishment is a trustworthy experimental substrate and a
locally verified next experiment. That is meaningful progress in research
method and AI-mathematics infrastructure. It is not evidence that AION is
closer to proving RH in the mathematical sense.

The completed v3 and retention results determine the architectural conclusion:

\begin{enumerate}
  \item v3 produced a directional mechanism signal (3/6 for full AION versus
  0/6 proposer), but missed the absolute 4/6 superiority floor; and
  \item delayed source-closed retention scored 2/6 for the memory-and-repair arm
  versus 3/6 for isolated cold AION, so it supplied no retained advantage.
\end{enumerate}

The programme therefore makes no AION-superiority claim. The governed
experiments remain useful negative evidence, while mathematical continuation
is assessed only through Lean-checked artifacts. That continuation has now
completed the critical-strip foundation, compact multiplicity-aware divisor,
scalar Perron theorem, and finite Chebyshev inversion.

\section{Post-seal execution addendum}

The fresh v3 matched comparison completed with zero infrastructure failures.
The primary scores were: proposer alone 0/6, full AION memory plus repair 3/6,
AION without memory 0/6, AION without repair 2/6, and cold AION with isolated
state 1/6. Full AION exceeded every comparator, but the precommitted absolute
floor was 4/6. Consequently the superiority gate was not met and no superiority
claim is supported.

The result is nevertheless diagnostically useful. Verified memory plus repair
recovered a scaled special value, a specialized trivial zero, and real-axis
positivity. Memory without repair recovered two tasks; deterministic repair
without memory recovered none; isolated cold AION recovered one. This is a
mechanism signal to replicate, not a scientific conclusion about general
capability.

The local formal foundation now contains 17 kernel-checked established results.
In particular, the programme reconstructs the strict upper boundary
\code{riemannZeta s = 0 -> s.re < 1} from Mathlib's nonvanishing theorem on the
closed half-plane. The remaining critical-strip blocker is precise: classify
every zero with nonpositive real part as a member of the negative-even trivial
zero family. This reconstruction is the next mathematical target; it is not a
claim of new mathematics.

The curriculum contains 36 governed tasks. The separate six-task retention
packet was sealed, executed exactly once after its recorded time gate, scored,
and closed. Its automation is disabled.

\begin{truthbox}
As of 15 August 2026: trusted formal laboratory - yes; governed memory
retrieval - yes; deterministic Mathlib discovery - yes; v3 and retention
experiments complete - yes; precommitted AION superiority - no; scalar Perron
theorem kernel-checked - yes; finite Chebyshev inversion kernel-checked - yes;
full Chebyshev Perron formula - yes; initial contour and residue ledger - yes;
global explicit formula - not yet; new original RH mathematics - none;
Riemann Hypothesis solved - no.
\end{truthbox}

\clearpage
\section{Mathematics continuation: compact zeta divisors and the explicit-formula interface}
\label{sec:compact-divisor-explicit-formula}

This section records the mathematical continuation completed after the v3
experiment. It supersedes the earlier forward-looking statements that named
left-half-plane classification or compact divisor construction as future work;
those established prerequisites have now been reconstructed in Lean. The
section is deliberately explicit about the boundary between what the kernel has
accepted and what remains an open analytic proof obligation.

The implementation recorded here is repository commit
\code{651f65f8}. The complete Lean library builds successfully against Lean
4.24.0 and pinned Mathlib v4.32.1. The research modules contain 55
kernel-checked theorems, or 58 declarations when the three baseline
reproduction checks are included. No project-authored \code{sorry},
\code{admit}, axiom, or unsafe declaration is used in the new modules.

\begin{truthbox}
This continuation formalizes established zeta-function infrastructure and
states the next explicit-formula theorem precisely. It does not prove the
explicit formula, improve a zero-free region, constrain a previously unknown
zero, or establish the Riemann Hypothesis. Its classification is
\emph{formalization progress}, not new mathematical progress on RH.
\end{truthbox}

\subsection{Nontrivial zeros and the reconstructed critical strip}

The local formalization distinguishes the negative-even trivial zeros from all
other zeta zeros. In mathematical notation,
\begin{align*}
  \operatorname{Triv}(s)
    &\;:\Longleftrightarrow\;
      \exists n\in\mathbb N,\quad s=-2(n+1),\\
  \operatorname{Nontriv}(s)
    &\;:\Longleftrightarrow\;
      \zeta(s)=0\quad\text{and}\quad\neg\operatorname{Triv}(s).
\end{align*}

The kernel-checked classification theorem reconstructs the established
left-half-plane result
\[
  \zeta(s)=0\ \wedge\ \RePart(s)\leq 0
  \quad\Longrightarrow\quad
  \operatorname{Triv}(s).
\]
Together with nonvanishing on the closed right half-plane, this yields
\[
  \operatorname{Nontriv}(s)
  \quad\Longrightarrow\quad
  0<\RePart(s)<1.
\]
The proof uses the completed functional equation, nonvanishing of the reflected
point $1-s$, and the exact zero classification of Mathlib's real Gamma factor.
It therefore reconstructs classical mathematics rather than assuming the
critical strip as a local axiom.

The standard symmetries are also checked:
\[
  \rho\longmapsto 1-\rho,
  \qquad
  \rho\longmapsto \overline{\rho},
\]
and hence every nontrivial zero generates the formal four-point orbit
\[
  \left\{\rho,\ 1-\rho,\ \overline{\rho},\
  1-\overline{\rho}\right\}.
\]
The local predicate is proved equivalent to Mathlib's formal statement of
\code{RiemannHypothesis}.

\subsection{Selected Chebyshev--psi criterion}

The programme now follows an RH-equivalent Chebyshev error-bound route. Define
the remainder
\[
  E_{\psi}(x):=\psi(x)-x,
\]
where
\[
  \psi(x)=\sum_{n\leq x}\Lambda(n)
\]
and $\Lambda$ is the von Mangoldt function. The selected square-root-scale
target is
\begin{equation}
  \exists C>0\ \ \forall x\geq 2,\qquad
  |E_{\psi}(x)|
  \leq C\sqrt{x}\,(\log x)^2.
  \label{eq:critical-psi-bound}
\end{equation}

The Lean development currently verifies the finite von Mangoldt representation
of $E_{\psi}$ and the right-half-plane logarithmic-derivative identity
\begin{equation}
  \sum_{n=1}^{\infty}\frac{\Lambda(n)}{n^s}
  =-\frac{\zeta'(s)}{\zeta(s)},
  \qquad \RePart(s)>1.
  \label{eq:von-mangoldt-log-derivative}
\end{equation}
It also records the equivalence between RH and
Equation~\eqref{eq:critical-psi-bound} as a proposition to reconstruct, not as
an axiom or completed theorem.

\subsection{Multiplicity as analytic vanishing order}

For $s\neq 1$, zeta is analytic at $s$. Its zero multiplicity is represented by
the analytic order
\[
  m(s):=\operatorname{ord}_{s}(\zeta)\in\mathbb N\cup\{\infty\}.
\]
The formal development proves:
\begin{align*}
  \operatorname{Nontriv}(\rho)
    &\Longrightarrow m(\rho)\neq 0,\\
  s\neq 1
    &\Longrightarrow
      \bigl(m(s)=0\Longleftrightarrow\zeta(s)\neq 0\bigr),\\
  s\neq 1
    &\Longrightarrow m(s)\neq\infty.
\end{align*}
The last statement is important: it justifies converting the extended natural
analytic order into an ordinary finite integer multiplicity. Its Lean proof
propagates finite order from $s=2$, where zeta is nonzero, across the connected
punctured plane $\mathbb C\setminus\{1\}$.

\subsection{Compact, multiplicity-aware zeta divisor}

Let $K\subset\mathbb C$ avoid the pole:
\[
  K\subseteq\mathbb C\setminus\{1\}.
\]
The local zeta divisor is the integer-valued locally finite function
\[
  D_K(s):=\operatorname{div}_{K}(\zeta)(s).
\]
For $s\in K$, the divisor value agrees with the finite analytic vanishing
order. Outside $K$ it is zero. The following properties are kernel-checked:
\begin{enumerate}
  \item if $K$ is compact, then
    $\operatorname{supp}(D_K)$ is finite;
  \item $D_K(s)\geq 0$ for every $s$, because zeta is analytic on $K$;
  \item for $s\in K$,
    \[
      s\in\operatorname{supp}(D_K)
      \quad\Longleftrightarrow\quad
      \zeta(s)=0;
    \]
  \item every selected nontrivial zero has strictly positive integer
    multiplicity; and
  \item for any weight $w:\mathbb C\to\mathbb C$, the multiplicity-aware sum
    \[
      \sum_{\rho\in\operatorname{supp}(D_K)}D_K(\rho)w(\rho)
    \]
    is a genuine finite sum.
\end{enumerate}

Pinned Mathlib exposes zeta as analytic on
$\mathbb C\setminus\{1\}$ but does not currently expose, through the imports
used by this programme, a global meromorphic-continuation theorem suitable for
silently crossing the pole. The explicit pole-avoidance hypothesis is therefore
a formal safety condition, not merely a presentation choice.

Representative Lean declarations are:
\begin{lstlisting}
noncomputable def zetaDivisorOn (K : Set C) :
    Function.locallyFinsuppWithin K Z :=
  divisor riemannZeta K

theorem zetaDivisorOn_support_finite {K : Set C}
    (hK : IsCompact K) :
    (zetaDivisorOn K).support.Finite

theorem mem_compactZetaDivisorSupport_iff_zero
    {K : Set C} (hK : IsCompact K)
    (havoid : K subset {1}^c) {s : C} (hs : s in K) :
    s in compactZetaDivisorSupport K hK <-> riemannZeta s = 0
\end{lstlisting}
As elsewhere in this report, ASCII placeholders in listings stand for the
corresponding Unicode Lean symbols used in the checked source.

\subsection{Pole-free compact exhaustion}

For $n\in\mathbb N$, define
\begin{equation}
  K_n:=\overline{B}(0,n)
       \setminus B\!\left(1,\frac{1}{n+1}\right).
  \label{eq:zeta-truncation-region}
\end{equation}
Each $K_n$ is compact and excludes $1$. The shrinking deleted ball is necessary
because a plain expanding closed ball eventually contains the pole, where the
available analytic divisor theorem cannot be applied.

The formalization proves the exhaustion property
\begin{equation}
  s\neq 1
  \quad\Longrightarrow\quad
  \exists n\in\mathbb N,\ s\in K_n.
  \label{eq:truncation-exhaustion}
\end{equation}
Filtering the finite divisor support by the nontrivial-zero predicate gives
\[
  \mathcal Z_n
  :=\left\{\rho\in\operatorname{supp}(D_{K_n}):
      \operatorname{Nontriv}(\rho)\right\}.
\]
The sets $\mathcal Z_n$ are finite, and every nontrivial zero belongs to at
least one of them. Thus the formal truncation cannot permanently omit an
arbitrary fixed nontrivial zero.

\subsection{Finite zero contribution and the exact open target}

The zero kernel required by the classical explicit formula is
\[
  W_x(\rho):=\frac{x^{\rho}}{\rho}.
\]
Using the compact divisor, the checked finite contribution is
\begin{equation}
  Z_n(x)
  :=\sum_{\rho\in\mathcal Z_n}
      D_{K_n}(\rho)\frac{x^{\rho}}{\rho}.
  \label{eq:finite-zero-contribution}
\end{equation}
This construction is multiplicity-aware by definition; repeated zeros do not
need to be represented by duplicated list entries.

For the current convention and nonintegral $x>1$, define the elementary
pole/trivial-zero correction
\begin{equation}
  A(x):=\log(2\pi)
  +\frac12\log\!\left(1-x^{-2}\right).
  \label{eq:archimedean-correction}
\end{equation}
The next analytic theorem has now been recorded exactly as
\code{ChebyshevExplicitFormulaTarget}:
\begin{equation}
  \forall x>1,\quad x\notin\mathbb N
  \quad\Longrightarrow\quad
  Z_n(x)\longrightarrow -E_{\psi}(x)-A(x).
  \label{eq:explicit-formula-target}
\end{equation}
Equivalently, once convergence and conventions are justified,
\[
  E_{\psi}(x)
  =-\lim_{n\to\infty}Z_n(x)-A(x).
\]
The nonintegrality hypothesis avoids selecting a left-, right-, or half-weight
value at jump points of $\psi$. A later extension may introduce the standard
symmetrized convention at prime-power jumps.

\begin{warningbox}
Equation~\eqref{eq:explicit-formula-target} is a typed Lean proposition, not a
proved theorem. Defining $Z_n(x)$ and proving that its summands are genuine
zeros does not establish convergence, justify a contour shift, or calculate
the required residues. Those are the current mathematical proof obligations.
\end{warningbox}

\clearpage
\subsection{Completed specialized Perron interface}
\label{subsec:perron-interface}

The first Perron unit has now been constructed in
\path{RiemannLab/Tasks/PerronFormula.lean}. Define the zeta
logarithmic-derivative kernel
\begin{equation}
  \mathcal K_{\zeta}(x,s)
  :=-\frac{\zeta'(s)}{\zeta(s)}\frac{x^s}{s}
  \label{eq:zeta-perron-kernel}
\end{equation}
and the equivalent von Mangoldt L-series kernel
\begin{equation}
  \mathcal K_{\Lambda}(x,s)
  :=\left(\sum_{n=1}^{\infty}\frac{\Lambda(n)}{n^s}\right)
    \frac{x^s}{s}.
  \label{eq:von-mangoldt-perron-kernel}
\end{equation}
For $\RePart(s)>1$, the previously verified logarithmic-derivative identity
gives the kernel-checked equality
\begin{equation}
  \mathcal K_{\Lambda}(x,s)=\mathcal K_{\zeta}(x,s).
  \label{eq:perron-kernel-equality}
\end{equation}

For $c>1$ and finite height $T$, parametrize the vertical segment by
$s=c+it$. Since $ds=i\,dt$, the factor $i$ cancels against the $i$ in
$(2\pi i)^{-1}$. The formal normalized truncations are therefore
\begin{align}
  P^{\zeta}_T(c,x)
    &:=\frac{1}{2\pi}\int_{-T}^{T}
      \mathcal K_{\zeta}(x,c+it)\,dt,
      \label{eq:truncated-zeta-perron}\\
  P^{\Lambda}_T(c,x)
    &:=\frac{1}{2\pi}\int_{-T}^{T}
      \mathcal K_{\Lambda}(x,c+it)\,dt.
      \label{eq:truncated-von-mangoldt-perron}
\end{align}
Lean proves that every point $c+it$ remains in the half-plane of absolute
convergence and hence, for every finite $T$,
\begin{equation}
  c>1
  \quad\Longrightarrow\quad
  P^{\Lambda}_T(c,x)=P^{\zeta}_T(c,x).
  \label{eq:truncated-perron-equality}
\end{equation}
The zero-height specializations are also checked:
\[
  P^{\zeta}_0(c,x)=P^{\Lambda}_0(c,x)=0.
\]

Representative Lean declarations are:
\begin{lstlisting}
noncomputable def zetaPerronKernel (x : R) (s : C) : C :=
  (-deriv riemannZeta s / riemannZeta s) * (x : C) ^ s / s

noncomputable def truncatedZetaPerronIntegral
    (c T x : R) : C :=
  (1 / (2 * (Real.pi : C))) *
    integral t in -T..T, zetaPerronKernel x (c + t * Complex.I)

theorem truncatedVonMangoldtPerronIntegral_eq_zeta
    {c T x : R} (hc : 1 < c) :
    truncatedVonMangoldtPerronIntegral c T x =
      truncatedZetaPerronIntegral c T x
\end{lstlisting}
The listing uses portable ASCII placeholders; the checked source uses Lean's
native Unicode notation and interval integral syntax.

\subsubsection{Exact inversion target}

For nonintegral $x>1$, the classical Perron conclusion required by the programme
is
\begin{equation}
  \forall c>1,\qquad
  \lim_{T\to\infty}P^{\zeta}_T(c,x)=\psi(x).
  \label{eq:chebyshev-perron-target}
\end{equation}
This is represented by the Lean proposition
\code{ChebyshevPerronInversionTarget}, and is now proved by
\code{chebyshevPerronInversion}. The proof passes through the arithmetic and
zeta forms using Equation~\eqref{eq:truncated-perron-equality}.

\begin{warningbox}
Perron inversion is complete away from jumps. The theorem deliberately assumes
$x$ is not a natural number, so no half-weight convention at a jump of
$\psi$ is being smuggled into the result.
\end{warningbox}

\subsubsection{Pinned-library capability audit}

The dependency audit found the following usable primitives in pinned Mathlib:
\begin{center}
\begin{tabularx}{\textwidth}{>{\raggedright\arraybackslash}p{55mm}X}
\toprule
Module (prefix \code{Mathlib.}) & Available role \\
\midrule
\path{Analysis.MellinTransform} &
Mellin transform, inverse Mellin transform, and vertical integrability. \\
\path{Analysis.MellinInversion} &
General Mellin inversion through Fourier inversion under its stated
integrability and continuity hypotheses. \\
\path{Analysis.Complex.CauchyIntegral} &
Rectangle boundary integrals and Cauchy--Goursat machinery needed for later
contour shifting. \\
\path{NumberTheory.LSeries.Dirichlet} &
Absolute convergence of the von Mangoldt L-series for $\RePart(s)>1$ and its
identity with $-\zeta'/\zeta$. \\
\bottomrule
\end{tabularx}
\end{center}

No specialized arithmetic Perron theorem matching
Equation~\eqref{eq:chebyshev-perron-target} was found. This negative result is
recorded in \path{perron_dependency_audit_v1.json}; it prevents the programme
from treating a general Mellin inversion theorem as though all hypotheses for
the discontinuous arithmetic step function had already been discharged.

\subsection{Remaining analytic bridge}

The Perron interface and the finite-height arithmetic/zeta identification are
now complete. The shortest credible path from this interface to
Equation~\eqref{eq:explicit-formula-target} is decomposed as follows.

\begin{enumerate}
  \item \textbf{Scalar Perron kernel limit.} For $c>0$, $y>0$, and $y\neq1$,
  prove
  \begin{equation}
    Q_T(c,y):=\frac{1}{2\pi}\int_{-T}^{T}
      \frac{y^{c+it}}{c+it}\,dt
    \longrightarrow
    \begin{cases}
      1,&y>1,\\
      0,&0<y<1.
    \end{cases}
    \label{eq:scalar-perron-target}
  \end{equation}
  The omitted point $y=1$ is the half-weight discontinuity. This isolates the
  Fourier/Mellin convergence issue from all arithmetic-series bookkeeping.

  \item \textbf{Finite Dirichlet-polynomial inversion.} Apply
  Equation~\eqref{eq:scalar-perron-target} term by term to a finite sum
  \[
    \sum_{n\leq N}\frac{\Lambda(n)}{n^s},
  \]
  producing the corresponding finite partial sum of $\psi(x)$.

  \item \textbf{Absolute tail control.} Use absolute convergence on
  $\RePart(s)=c>1$ to control the omitted Dirichlet-series tail uniformly
  enough to pass from finite polynomials to
  Equation~\eqref{eq:chebyshev-perron-target}.

  \item \textbf{Rectangular contour shift.} Move the vertical line from
  $\RePart(s)=c$ into the left half-plane while avoiding zeros and the pole.
  Formal obligations include path composition, orientation, integrability,
  horizontal-edge bounds, and a sequence of admissible heights.

  \item \textbf{Residue accounting.} Prove that the pole at $s=1$ supplies the
  main term $x$, each nontrivial zero $\rho$ supplies
  $-m(\rho)x^{\rho}/\rho$, and the remaining pole/kernel and trivial-zero
  contributions combine into Equation~\eqref{eq:archimedean-correction}.

  \item \textbf{Limit and ordering control.} Relate the zeros encountered by
  rectangular contours to the compact exhaustion $K_n$, prove the required
  convergence mode, and justify the symmetric ordering of the conditionally
  convergent zero contribution.

  \item \textbf{RH-equivalent estimates.} Only after the explicit formula is
  established should the programme prove RH implies
  Equation~\eqref{eq:critical-psi-bound}, and conversely use the psi bound plus
  analytic continuation to exclude zeros with real part greater than $1/2$.
\end{enumerate}

Items 1--3 are complete. The irreducible unit is now the global punctured
rectangle theorem: excise finitely many poles, relate the small-circle
integrals to the checked local residues, and prove the displaced vertical and
horizontal edges vanish along admissible heights.

\begin{decisionbox}
The specialized Perron representation and its limit to $\psi(x)$ are checked.
The next work must complete genuine pole excision and boundary estimates; a
formal residue ledger alone is not a contour-shift theorem.
\end{decisionbox}

\subsection{Sharp convergence and absolute Fourier reduction completed}

The conditional-convergence obstacle has now been separated from the value
calculation in \path{RiemannLab/Tasks/OscillatoryCauchy.lean}. For $c>0$ and
$a\neq0$, define
\[
  H_T(c,a):=\frac{1}{2\pi}\int_{-T}^{T}
    \frac{e^{iat}}{c+it}\,dt.
\]
Lean verifies an exact finite-interval integration-by-parts identity, proves
that both endpoint terms vanish like $O((|a|T)^{-1})$, and proves absolute
integrability of the squared-denominator remainder. Consequently the sharp
symmetric truncations $H_T(c,a)$ genuinely converge; this is no longer an open
assumption.

The same module now proves the source-closed Laplace identity
\[
  \frac{1}{c+it}=\int_0^\infty e^{-(c+it)u}\,du
\]
and the exact reduction
\begin{equation}
  \lim_{T\to\infty}H_T(c,a)
  =\frac{1}{2\pi a}\int_{-\infty}^{\infty}
    \frac{e^{iat}}{(c+it)^2}\,dt.
  \label{eq:squared-cauchy-reduction}
\end{equation}
The integral on the right of
Equation~\eqref{eq:squared-cauchy-reduction} is absolutely convergent. Thus the
remaining analytic gate has the single named statement
\begin{equation}
  \int_{-\infty}^{\infty}\frac{e^{iat}}{(c+it)^2}\,dt
  =
  \begin{cases}
    2\pi a e^{-ac},&a>0,\\
    0,&a<0.
  \end{cases}
  \label{eq:squared-cauchy-value}
\end{equation}
Lean also verifies that Equation~\eqref{eq:squared-cauchy-value} implies the
original oscillatory Cauchy target and hence the away-from-jump scalar Perron
limit. The preferred next route is Fourier inversion for the continuous
one-sided function $u\mapsto u e^{-cu}\mathbf{1}_{(0,\infty)}(u)$, whose
transform has precisely the squared denominator in
Equation~\eqref{eq:squared-cauchy-value}.

\begin{resultbox}
At this intermediate checkpoint, sharp truncation, endpoint decay, and
conditional tail control reduced the scalar Perron gate to one absolutely
convergent Fourier value. The following section records the subsequent closure
of that value.
\end{resultbox}

\subsection{Fourier value and scalar Perron theorem completed}

The final scalar value has now been kernel-checked in
\path{RiemannLab/Tasks/SquaredCauchyFourier.lean}. For $c>0$, Lean constructs
the continuous one-sided exponential moment
\[
  f_c(u)=\begin{cases}u e^{-cu},&u>0,\\0,&u\leq0,\end{cases}
\]
proves its whole-line integrability, computes its Fourier transform, verifies
the integrability of that transform, and applies the pinned Mathlib Fourier
inversion theorem. This proves Equation~\eqref{eq:squared-cauchy-value} and
propagates through the prior reductions to establish
\[
  \frac1{2\pi}\int_{-T}^{T}\frac{y^{c+it}}{c+it}\,dt
  \longrightarrow
  \begin{cases}1,&y>1,\\0,&0<y<1,\end{cases}
\]
for every $c>0$ away from the jump $y=1$. The scalar Perron target is therefore
no longer open.

\subsection{Finite von Mangoldt and Chebyshev inversion completed}

The new module \path{RiemannLab/Tasks/FinitePerron.lean} lifts the scalar
result to an arbitrary finite coefficient family. In particular, for
nonintegral $x>0$ it proves the finite von Mangoldt specialization
\begin{equation}
  \sum_{n\leq\lfloor x\rfloor}\Lambda(n)
    \frac1{2\pi}\int_{-T}^{T}
      \frac{(x/n)^{c+it}}{c+it}\,dt
  \longrightarrow \psi(x).
  \label{eq:finite-chebyshev-perron}
\end{equation}
All sums in Equation~\eqref{eq:finite-chebyshev-perron} are finite, so this
theorem contains no hidden exchange of an infinite series and an integral.

The full Perron module also proves fixed-height interchange of the absolutely
convergent von Mangoldt series and the vertical integral. A
height-independent Dirichlet bound supplies a summable coefficient majorant,
so Tannery's theorem passes the growing-height limit through the series. Thus
the complete zeta Perron integral converges to $\psi(x)$ away from integer
jumps.

\begin{resultbox}
The scalar theorem, finite inversion, fixed-height interchange, uniform
coefficient domination, and the complete Chebyshev Perron formula are all
kernel-checked. This is significant formalization of classical analysis, not
a new theorem about zeta zeros and not a proof of RH.
\end{resultbox}

\subsection{Formal artifact map}

\begin{longtable}{@{}p{47mm}p{101mm}@{}}
\toprule
Artifact & Technical role \\
\midrule
\endhead
\path{RiemannLab/Tasks/NontrivialZeros.lean} &
Trivial/nontrivial predicates, critical-strip classification, reflection,
conjugation, four-point orbit, and RH equivalence. \\
\path{RiemannLab/Tasks/CriticalLineCriterion.lean} &
Chebyshev remainder, square-root-scale criterion, finite von Mangoldt sum, and
zeta logarithmic-derivative bridge. \\
\path{RiemannLab/Tasks/ZetaZeroMultiplicity.lean} &
Analytic zero order, analyticity at nontrivial zeros, and zero/nonzero
characterization away from the pole. \\
\path{RiemannLab/Tasks/CompactZetaDivisor.lean} &
Locally finite integer divisor, compact support, finite order, support/zero
equivalence, nonnegativity, and finite weighted sums. \\
\path{RiemannLab/Tasks/ExplicitFormula.lean} &
Pole-free compact exhaustion, finite nontrivial-zero sets, positive
multiplicities, finite $x^\rho/\rho$ sums, and the open convergence target. \\
\path{RiemannLab/Tasks/PerronFormula.lean} &
Normalized zeta and von Mangoldt Perron kernels, finite-height vertical
integrals, exact equality on $\RePart(s)>1$, fixed-height interchange,
summable uniform majorant, Tannery limit, and full Chebyshev Perron formula. \\
\path{RiemannLab/Tasks/ContourShift.lean} &
Oriented rectangle boundary, normalized right-edge identity, local
multiplicity-aware zero residues, and residues at $s=1$ and $s=0$. \\
\path{RiemannLab/Tasks/ScalarPerron.lean} &
Scalar kernel, logarithmic reduction, half-weight value at the jump, and exact
equivalence with the oscillatory Cauchy limit. \\
\path{RiemannLab/Tasks/OscillatoryCauchy.lean} &
Sharp convergence, integration-by-parts tail control, Laplace representation,
and the absolute squared-kernel reduction subsequently discharged by Fourier
inversion. \\
\path{RiemannLab/Tasks/SquaredCauchyFourier.lean} &
One-sided exponential moment, Fourier transform and inversion, exact
squared-Cauchy value, and completed scalar Perron theorem. \\
\path{RiemannLab/Tasks/FinitePerron.lean} &
Generic finite coefficient inversion and finite von Mangoldt specialization
to the Chebyshev function away from integer jumps. \\
\path{explicit_formula_dependency_graph_v1.json} &
Machine-readable dependency state and the precommitted next proof unit. \\
\path{perron_dependency_audit_v1.json} &
Pinned Mathlib capability audit, missing specialized theorem, and scalar
Perron-kernel decomposition. \\
\bottomrule
\end{longtable}

\begin{resultbox}
The compact divisor and complete Chebyshev Perron formula are no longer
blockers. Local zero residues with full multiplicity, the residues at $s=1$
and $s=0$, rectangle orientation, and the right-edge normalization are also
verified. The remaining classical blocker is global pole excision together
with trivial-zero summation and vanishing-edge estimates. After the explicit
formula, the programme reaches the exact classical RH frontier.
\end{resultbox}

\section{Full Perron inversion and initial contour residues completed}

The module \path{RiemannLab/Tasks/PerronFormula.lean} now proves
\begin{equation}
  \lim_{T\to\infty}\frac{1}{2\pi}
  \int_{-T}^{T}-\frac{\zeta'(c+it)}{\zeta(c+it)}
    \frac{x^{c+it}}{c+it}\,dt=\psi(x)
\end{equation}
for $c>1$, $x>1$, and nonintegral $x$. The decisive new estimate bounds each
normalized coefficient integral by a summable majorant independent of $T$;
Tannery's theorem then passes the limit through the infinite von Mangoldt
series.

The new module \path{RiemannLab/Tasks/ContourShift.lean} fixes the four-side
orientation of the rectangle and proves that its normalized right edge is
exactly the established finite-height Perron integral. Analytic factorization
at a zero gives the multiplicity-aware limit
\begin{equation}
  (s-\rho)\frac{\zeta'(s)}{\zeta(s)}\longrightarrow m(\rho),
\end{equation}
and therefore the contour kernel has local residue
$-m(\rho)x^\rho/\rho$. The compact divisor converts these local values into a
finite residue ledger with the same multiplicities used by the truncated zero
sum. Mathlib's Laurent asymptotic at $s=1$ yields the main term $x$, while the
known values $\zeta(0)=-1/2$ and
$\zeta'(0)=-\log(2\pi)/2$ yield the residue $-\log(2\pi)$ at the Perron pole.

These are classical results formalized with kernel checking. They do not yet
constitute the explicit formula: admissible-height estimates for the remaining
three contour edges remain open.

\section{Global residue theorem, trivial-zero sum, and final assembly}

The classical contour infrastructure has advanced through three further
machine-checked gates.  The module
\path{RiemannLab/Tasks/TrivialZeroResidues.lean} proves that every negative-even
zero $-2(n+1)$ is simple.  It identifies its contour residue as
\[
  -\frac{x^{-2(n+1)}}{-2(n+1)}
  = \frac{x^{-2(n+1)}}{2(n+1)}
\]
and evaluates the full series as
\[
  \sum_{n\geq0}\frac{x^{-2(n+1)}}{2(n+1)}
  =-\frac12\log(1-x^{-2}), \qquad x>1.
\]
Thus the trivial-zero simplicity and summation gate is complete.

The module \path{RiemannLab/Tasks/RectangularResidues.lean} proves the missing
rectangular Cauchy kernel identity directly in the exact four-edge orientation
used by the project:
\[
  \int_{\partial R}\frac{ds}{s-p}=2\pi i
  \qquad (p\in\operatorname{int}R).
\]
The proof controls the branch cut by using $\log(s-p)$ on three sides and
$\log(-(s-p))$ on the left side.  It then lifts this identity to a finite
global residue theorem.  For a holomorphic remainder $g$, a finite pole set
$P$, pole locations $p_j$, and residues $r_j$, Lean proves
\[
 \int_{\partial R}
 \left(g(s)+\sum_{j\in P}\frac{r_j}{s-p_j}\right)ds
 =2\pi i\sum_{j\in P}r_j,
\]
provided every $p_j$ lies strictly inside $R$.  Continuity supplies the edge
integrability automatically.  This closes the global finite residue-collection
gate without assuming a residue theorem absent from the pinned Mathlib version.

Finally, \path{RiemannLab/Tasks/ExplicitFormulaAssembly.lean} defines the exact
remaining admissible-contour limit proposition and proves that it implies
\code{ChebyshevExplicitFormulaTarget}.  In particular, once a sequence of
admissible heights tends to infinity, the normalized top, bottom, and displaced
left edges tend to zero, and the finite residue ledger is instantiated, no
further complex analysis or limit algebra remains: the already proved Perron
limit forces the complete Chebyshev explicit formula.

The full pinned build at that checkpoint passed in 3,593 jobs.  Two tightly coupled obligations
remain: instantiate the finite principal-parts theorem for the actual zeta
kernel on each chosen rectangle, and prove the zeta logarithmic-derivative
estimate needed to construct admissible heights and make the three displaced
edge integrals vanish.  The first is now bounded analytic packaging around the
proved residue theorem; the second is the substantive remaining classical
estimate. Pinned Mathlib supplies neither that growth estimate nor a ready
theorem from which it can be obtained. This is still established classical
analytic number theory being formalized; it is not new evidence for RH and
imposes no new constraint on the real parts of nontrivial zeros.

\section{Finite principal parts for the actual zeta kernel}

The analytic-packaging obligation from the preceding section is now complete.
The new module \path{RiemannLab/Tasks/MeromorphicPrincipalPart.lean} proves a
generic removal theorem.  If $F$ is meromorphic at $p$ and
\[
  (s-p)F(s)\longrightarrow r
  \qquad(s\to p,\ s\ne p),
\]
then there is a function analytic at $p$ which agrees punctured-locally with
$F(s)-r/(s-p)$.  The proof is carried out through Mathlib's meromorphic order:
multiplication by $(s-p)$ raises the order by one, and convergence to zero
forces the remaining order to be nonnegative.

The module \path{RiemannLab/Tasks/FinitePrincipalParts.lean} lifts this local
statement to an arbitrary finite pole set $P$.  Given residues $r_p$, it
constructs a single function $g$ analytic at every $s\in\mathbb C$ such that
\[
  F(s)=g(s)+\sum_{p\in P}\frac{r_p}{s-p},\qquad s\notin P.
\]
The construction patches the removable values at the finitely many poles and
proves analyticity both at each patched point and on the complement.  Thus the
finite residue theorem can now be applied without assuming that independently
chosen local extensions form a global holomorphic remainder.

The module \path{RiemannLab/Tasks/ZetaContourPrincipalParts.lean} instantiates
this theorem for the actual contour kernel
\[
  K_x(s)=-\frac{\zeta'(s)}{\zeta(s)}\frac{x^s}{s}.
\]
It proves global meromorphy of $\zeta$ from its completed-zeta representation
at $s=1$ and analyticity elsewhere, hence global meromorphy of $K_x$.  The
principal parts are then removed at all four pole classes:
\[
\begin{array}{c|c}
  \text{point} & \operatorname*{Res} K_x \\
  \hline
  1 & x \\
  0 & -\log(2\pi) \\
  -2(n+1) & -x^{-2(n+1)}/[-2(n+1)] \\
  \rho\ \text{nontrivial} & -m(\rho)x^\rho/\rho.
\end{array}
\]
For any finite ledger whose entries carry one of these classifications, and
which contains every singularity under discussion, Lean now produces the
globally analytic remainder required by the rectangular residue theorem.

\section{Quantitative reduction of the contour edges}

The new module \path{RiemannLab/Tasks/ContourEdgeEstimates.lean} isolates the
remaining large-height input by proving the exact identity
\[
  \lVert K_x(s)\rVert
  =\left\lVert\frac{\zeta'(s)}{\zeta(s)}\right\rVert
    \frac{x^{\Re s}}{\lVert s\rVert},\qquad x>0.
\]
It then converts any supplied bound
$\lVert\zeta'(s)/\zeta(s)\rVert\leq L$ and lower bound
$R\leq\lVert s\rVert$ into
\[
  \lVert K_x(s)\rVert\leq \frac{Lx^{\Re s}}{R}.
\]
Uniform pointwise bounds are lifted to the exact horizontal and vertical
interval integrals used by the project.  In particular, if $C_n$ uniformly
bounds a horizontal edge from $a$ to $b$ and
$C_n|b-a|\to0$, Lean proves that the corresponding complex edge integral
tends to zero.

The full pinned build now passes in 3,599 jobs.  The finite principal-parts
instantiation is therefore closed.  What remains is the substantive classical
large-height theorem: construct heights avoiding zero ordinates with a uniform
bound for $\zeta'/\zeta$, discharge the top and bottom edges through the new
norm estimates, and combine the functional equation and gamma-factor bounds
with a displaced-left sequence to make the left edge vanish.  The finite
ledger must finally be identified with the compact multiplicity-aware zero
truncation already used by the target statement.

Pinned Mathlib does not contain this global admissible-height estimate.  No
complete Chebyshev explicit formula is claimed at this checkpoint.  All new
results remain formalizations of classical analytic infrastructure and provide
no new restriction on nontrivial zero locations.

\section{Geometric admissible heights and exact three-edge assembly}

The geometric portion of the admissible-height construction is now
machine-checked.  The module
\path{RiemannLab/Tasks/AdmissibleHeights.lean} selects, for every
$n\in\mathbb N$, a height
\[
  n+2<T_n<n+3
\]
such that
\[
  \zeta(\sigma+iT_n)\ne0\qquad(\sigma\in\mathbb R).
\]
Compact discreteness supplies only finitely many zeros in the relevant closed
ball.  An ordinate outside that finite set is then selected inside the unit
interval.  The classification of zeros with nonpositive real part and the
critical-strip location theorem ensure that every possible zero on the line
would have appeared in that compact set.  Complex conjugation gives the
reflected zero-free line at $-T_n$, and Lean proves $T_n\to\infty$.

On every fixed compact interval $[a,b]$, continuity of $\zeta'/\zeta$ on the
two selected lines now yields a finite common uniform bound.  This deliberately
does not assert a rate.  The exact quantitative proposition isolated in
\path{RiemannLab/Tasks/AdmissibleHeightBounds.lean} requires bounds $B_n$ with
\[
  \sup_{-1\leq\sigma\leq c}
  \left|\frac{\zeta'}{\zeta}(\sigma\mathbin{\pm}iT_n)\right|
  \leq B_n,
  \qquad \frac{B_n}{T_n}\longrightarrow0.
\]
From this input Lean proves that the corresponding upper and lower
fixed-strip kernel integrals vanish.

The module \path{RiemannLab/Tasks/DisplacedLeftEdge.lean} defines the
displaced real parts
\[
  L_n=-(2n+1).
\]
These lines tend to $-\infty$ and avoid every zeta zero: nontrivial zeros lie
in $0<\Re s<1$, whereas trivial zeros are negative even integers.  The file
states the exact remaining functional-equation and Gamma-factor decay target
and proves that it forces the oriented left-edge integral to vanish.

Finally, \path{RiemannLab/Tasks/AdmissibleContourEdges.lean} assembles the
bottom, top, and displaced-left sides with the precise signs inherited from
the verified rectangle orientation.  For
\[
  K_x(s)=-\frac{\zeta'(s)}{\zeta(s)}\frac{x^s}{s},
\]
the normalized error is
\[
 E_n(x,c)=\frac{1}{2\pi i}\left(
   \int_{L_n}^{c}K_x(\sigma-iT_n)\,d\sigma
  -\int_{L_n}^{c}K_x(\sigma+iT_n)\,d\sigma
  -i\int_{-T_n}^{T_n}K_x(L_n+it)\,dt\right).
\]
Lean proves $E_n(x,c)\to0$ once the two expanding horizontal integrals and the
left integral vanish.  Thus no sign or normalization step remains hidden in
the final contour assembly.

The full pinned build passes in \textbf{3,603 jobs}.  The unresolved theorem is
now stated more accurately than at the previous checkpoint.  Merely avoiding
zero ordinates is insufficient: one needs quantitative separation from nearby
zeros together with a zero-counting/logarithmic-derivative estimate strong
enough on the expanding interval $[L_n,c]$.  One also needs the displaced-left
decay from the functional equation and quantitative Gamma estimates.  The
finite rectangle ledger must then be matched to the compact nontrivial-zero
truncation and to the vanishing tail of the already summed trivial-zero series.
These are established classical results, but they are not present globally in
the pinned Mathlib library and have not yet been proved in this project.
Consequently neither the complete explicit formula nor the classical RH
frontier is claimed at this checkpoint, and no new progress on RH itself is
claimed.

\appendix

\section{Key evidence hashes}

\begin{longtable}{@{}p{52mm}p{96mm}@{}}
\toprule
Artifact & SHA-256 \\
\midrule
\endhead
Immutable paid live result & \hash{5795e3faa55cf2a7466a3537ac41005056d44711eaa826c84935af318437b2e6} \\
Established zeta baseline & \hash{d121313f104e8e9d3a62fd2270e46aa6c4a7c6888ebc7d996f004fae5a25943c} \\
Catalogue probe & \hash{6691c0dff28a5b761b7c4d1da63f7820cb0b57976b21979594be093e149fd5b5} \\
V2 packet & \hash{2190c93a840eebabf0ad7cbfc3dbfe53a01ca32ead4f56765143d190e50d3660} \\
Sealed v3 packet & \hash{1771f8b9977ce0cf1ffac73b1284361e61ae92d4f86672afe97a241a786abb6e} \\
Completed v3 result & \hash{52899a504530456f1c6f90acf28a34edddc54ccbb4ce765d145c2f9166bc951d} \\
Pinned v2 environment & \hash{6776200251cfd041ecda62038ae8ff02980bf1066a332289b6c96ca015485591} \\
\bottomrule
\end{longtable}

The local-verification report is regenerated and therefore receives a new hash
when timestamps or verified memory state are rebuilt. The timestamped paid
result is the immutable historical object.

\section{V2 statement hashes}

\begin{longtable}{@{}p{11mm}p{137mm}@{}}
\toprule
Ref. & Statement SHA-256 \\
\midrule
\endhead
T1 & \hash{f84686a87dff3c814d79a296232fad81531194b318c4c81fe7d5aab8c9050798} \\
T2 & \hash{d1ca22f154945c4c26cf5f5cecfc5bf86d52d1538375d94420ea5abf7fce2760} \\
T3 & \hash{3c6a7aa1eede91cbd24530d477c38f750176269d1d8c08773dd2bc510fe77d31} \\
T4 & \hash{37c2b2c0f357b924d945d3c815244f0eff7affc33372d60ba839caea750cace5} \\
T5 & \hash{fed50dcca60493a2beedc23cb7ea0d485c1659dbe0fabdd0b8d3e698db523d3a} \\
T6 & \hash{6e920be25538210c2ec4850febc0c8a50e426af5c4ba743253c84b8af3611b4b} \\
\bottomrule
\end{longtable}

\section{Reproducibility commands}

Run from the COMDEX repository root:

\begin{lstlisting}[language={}]
.venv/bin/python scripts/aion_riemann_local_verify.py
.venv/bin/python -m pytest backend/tests/test_aion_riemann_lab.py -q
\end{lstlisting}

The v3 live comparison and the delayed retention comparison are complete and
protected by their one-shot execution receipts. Neither result established an
AION advantage, and no further paid run is scheduled by this document.

\section{Selected references}

\begin{thebibliography}{9}
\bibitem{Riemann1859}
B. Riemann, ``On the Number of Prime Numbers Less Than a Given Quantity,'' 1859.

\bibitem{Titchmarsh}
E. C. Titchmarsh, revised by D. R. Heath-Brown, \emph{The Theory of the Riemann Zeta-Function}, Oxford University Press.

\bibitem{Edwards}
H. M. Edwards, \emph{Riemann's Zeta Function}, Dover Publications.

\bibitem{Lean}
Lean Prover community, \emph{Lean 4 theorem prover and reference documentation}.

\bibitem{Mathlib}
Mathlib community, \emph{Mathlib: a mathematical library for Lean}, pinned in this programme at v4.32.1.

\bibitem{AionEvidence}
AION/COMDEX, \emph{Riemann Hypothesis Programme continuation handover, local verification report, task packets, Lean sources, and immutable JSON evidence}, 14--15 August 2026.
\end{thebibliography}

\end{document}
